% GENERATED FILE. Edit canonical lessons, then run npm run generate:books.
% canonical-source-hash: fnv1a64:db0a8ddb3d4650cf
% canonical-lessons: HI-C236-adha, HI-C236-pura, HI-C236-dugna, HI-C236-alag, HI-C236-gol

\chapter{Half, Whole, Double, Separate, Round}
\label{ch:hi-half-whole-double-separate-round}
\hlchaptermodality{236}

\begin{chapteropening}
\textbf{By the end of this chapter:} \emph{I can say half, whole, double, separate and round.}

Complete the last lesson of chapter 236: i can say half, whole, double, separate and round.
\end{chapteropening}

\section[\texorpdfstring{\hi{आधा} (ādhā)}{आधा (ādhā)}]{\hi{आधा} (ādhā) — half}

\label{lesson:HI-C236-adha}



\begin{quote}
Before the new one: say the Hindi for behind, then the Hindi for inside.
\end{quote}

\subsection*{You'll want to know: \hi{आधा}}
\textbf{\hi{आधा}} — \emph{ādhā} — \textquotedblleft{}half\textquotedblright{}.

\begin{grammarlens}[title={What you've built}]
One more word for this chapter.
\end{grammarlens}

\subsection*{Your turn}
\begin{itemize}
\raggedright
  \item \emph{Say it:} \emph{ādhā}
  \item \emph{Say it:} \emph{ādhā}, once more
  \item \emph{Say it:} say \emph{ādhā}
  \item \emph{Recall:} say the Hindi for behind, then the Hindi for inside, then say \emph{ādhā} again
\end{itemize}

\begin{tcolorbox}[breakable,colback=teal!4,colframe=teal!35!black,title={Before you move on}]
What does \hi{आधा} mean? (\textquotedblleft{}Half\textquotedblright{}.) Say it once more.
\end{tcolorbox}
\section[\texorpdfstring{\hi{पूरा} (pūrā)}{पूरा (pūrā)}]{\hi{पूरा} (pūrā) — whole, complete}

\label{lesson:HI-C236-pura}



\begin{quote}
Before the new one: say the Hindi for inside, then the Hindi for half.
\end{quote}

\subsection*{You'll want to know: \hi{पूरा}}
\textbf{\hi{पूरा}} — \emph{pūrā} — \textquotedblleft{}whole, complete\textquotedblright{}.

\begin{grammarlens}[title={What you've built}]
One more word for this chapter.
\end{grammarlens}

\subsection*{Your turn}
\begin{itemize}
\raggedright
  \item \emph{Say it:} \emph{pūrā}
  \item \emph{Say it:} \emph{pūrā}, once more
  \item \emph{Say it:} say \emph{pūrā}
  \item \emph{Recall:} say the Hindi for inside, then the Hindi for half, then say \emph{pūrā} again
\end{itemize}

\begin{tcolorbox}[breakable,colback=teal!4,colframe=teal!35!black,title={Before you move on}]
What does \hi{पूरा} mean? (\textquotedblleft{}Whole, complete\textquotedblright{}.) Say it once more.
\end{tcolorbox}
\section[\texorpdfstring{\hi{दुगना} (dugnā)}{दुगना (dugnā)}]{\hi{दुगना} (dugnā) — double}

\label{lesson:HI-C236-dugna}



\begin{quote}
Before the new one: say the Hindi for half, then the Hindi for whole.
\end{quote}

\subsection*{You'll want to know: \hi{दुगना}}
\textbf{\hi{दुगना}} — \emph{dugnā} — \textquotedblleft{}double\textquotedblright{}.

\begin{grammarlens}[title={What you've built}]
One more word for this chapter.
\end{grammarlens}

\subsection*{Your turn}
\begin{itemize}
\raggedright
  \item \emph{Say it:} \emph{dugnā}
  \item \emph{Say it:} \emph{dugnā}, once more
  \item \emph{Say it:} say \emph{dugnā}
  \item \emph{Recall:} say the Hindi for half, then the Hindi for whole, then say \emph{dugnā} again
\end{itemize}

\begin{tcolorbox}[breakable,colback=teal!4,colframe=teal!35!black,title={Before you move on}]
What does \hi{दुगना} mean? (\textquotedblleft{}Double\textquotedblright{}.) Say it once more.
\end{tcolorbox}
\section[\texorpdfstring{\hi{अलग} (alag)}{अलग (alag)}]{\hi{अलग} (alag) — separate, different}

\label{lesson:HI-C236-alag}



\begin{quote}
Before the new one: say the Hindi for whole, then the Hindi for double.
\end{quote}

\subsection*{You'll want to know: \hi{अलग}}
\textbf{\hi{अलग}} — \emph{alag} — \textquotedblleft{}separate, different\textquotedblright{}.

\begin{grammarlens}[title={What you've built}]
One more word for this chapter.
\end{grammarlens}

\subsection*{Your turn}
\begin{itemize}
\raggedright
  \item \emph{Say it:} \emph{alag}
  \item \emph{Say it:} \emph{alag}, once more
  \item \emph{Say it:} say \emph{alag}
  \item \emph{Recall:} say the Hindi for whole, then the Hindi for double, then say \emph{alag} again
\end{itemize}

\begin{tcolorbox}[breakable,colback=teal!4,colframe=teal!35!black,title={Before you move on}]
What does \hi{अलग} mean? (\textquotedblleft{}Separate, different\textquotedblright{}.) Say it once more.
\end{tcolorbox}
\section[\texorpdfstring{\hi{गोल} (gol)}{गोल (gol)}]{\hi{गोल} (gol) — round}

\label{lesson:HI-C236-gol}



\begin{quote}
Before the new one: say the Hindi for double, then the Hindi for separate.
\end{quote}

\subsection*{You'll want to know: \hi{गोल}}
\textbf{\hi{गोल}} — \emph{gol} — \textquotedblleft{}round\textquotedblright{}.

\begin{grammarlens}[title={What you've built}]
One more word for this chapter.
\end{grammarlens}

\subsection*{Your turn}
\begin{itemize}
\raggedright
  \item \emph{Say it:} \emph{gol}
  \item \emph{Say it:} \emph{gol}, once more
  \item \emph{Say it:} say \emph{gol}
  \item \emph{Recall:} say the Hindi for double, then the Hindi for separate, then say \emph{gol} again
\end{itemize}

\begin{tcolorbox}[breakable,colback=teal!4,colframe=teal!35!black,title={Before you move on}]
What does \hi{गोल} mean? (\textquotedblleft{}Round\textquotedblright{}.) Say it once more.
\end{tcolorbox}
